\section{Evaluation}
\label{sec:eval}

\subsection{Experimental Setup}

We evaluate ADSD in the reproduced HSD workspace of
\citet{zhou2026hsd} using \texttt{Qwen2.5-0.5B-Instruct-GPTQ-Int8} as the
draft model and \texttt{Qwen2.5-14B-Instruct-GPTQ-Int8} as the target model
under token-wise speculative verification. Unless otherwise noted, the
benchmark is GSM8K with the reproduced Qwen chat template and a matched
five-token benign prefix budget. All end-to-end numbers are reported against
the same benign token-wise baseline.

Each run uses batch size one, stochastic decoding
(\texttt{do\_sample=True}), and a per-sample cap of
\texttt{max\_new\_tokens=512}. We measure systems impact with wall-clock
latency, aggregate acceptance, block efficiency, and decoding speed. Latency
is recorded directly around the speculative \texttt{generate(...)} call.
Aggregate acceptance is the fraction of drafted tokens that survive
verification, and block efficiency is the number of accepted drafted tokens per
target verifier step. These counter-based metrics are computed from the saved
speculative traces in the reproduced HSD evaluator and therefore expose the
behavior of the verification mechanism independently of incidental timing
noise. Evaluating at batch size one follows standard speculative-decoding
practice \citep{leviathan2023speculative,zhou2026hsd}: the object of study is
the per-sequence verifier mechanism itself rather than downstream serving-stack
scheduling.

Before the end-to-end benchmarks, the Palmetto search consistently finds
projected prefixes that collapse early verifier acceptance. In the single-prompt
setting, the optimized five-token prefix reduces exact projected
$\alpha_1$ from $0.1597$ to $0.0288$. In the calibration-batch universal search
with $K=64$, every calibration size $N\in\{1,2,4,8\}$ yields a projected
prefix with calibration-batch $\alpha_1=0$, and the $N=4$ prefix delivers the
strongest accuracy-preserving end-to-end attack.

\subsection{End-to-End Latency Degradation}

ADSD substantially degrades verifier efficiency on unseen prompts while
preserving task accuracy. Table~\ref{tab:tokenwise-benchmark} reports the
matched 263-sample GSM8K benchmark comparing the benign boundary baseline with
the strongest accuracy-preserving universal prefix ($N=4$) and the strongest
cost-maximizing universal prefix ($N=2$).

\begin{table}[t]
\centering
\caption{Token-wise 263-sample GSM8K benchmark. Lower aggregate acceptance,
block efficiency, and decoding speed are worse for the system; higher mean
sample time is worse for the system.}
\label{tab:tokenwise-benchmark}
\footnotesize
\begin{tabular}{lccccc}
\toprule
Setting & Time (s) & Agg. acc. & Block eff. & Speed & GSM8K acc. \\
\midrule
Benign baseline & $23.10$ & $0.597$ & $5.832$ & $79.55$ & $222/263$ \\
Universal prefix ($N=4$) & $38.61$ & $0.390$ & $3.731$ & $51.68$ & $222/263$ \\
Universal prefix ($N=2$) & $51.60$ & $0.358$ & $3.506$ & $46.09$ & $202/263$ \\
\bottomrule
\end{tabular}
\end{table}

The systems signal is unambiguous. The $N=4$ universal prefix increases mean
sample time from $23.10$s to $38.61$s, a $67.1\%$ slowdown, while lowering
aggregate acceptance from $0.597$ to $0.390$ and block efficiency from $5.832$
to $3.731$. The same condition reduces HSD-style decoding speed from $79.55$
to $51.68$ while preserving GSM8K accuracy exactly at $222/263$. This
alignment between counter-based verifier statistics and wall-clock latency is
critical: it shows that the slowdown is a consequence of algorithmic verifier
degradation rather than an artifact of hardware noise. The $N=2$ universal
prefix pushes the cost attack further still, but at the price of task
degradation. For the core stealth claim of the paper, the $N=4$ prefix is the
most important result: a short prompt-side perturbation can deny the intended
speculative speedup while leaving user-visible task quality unchanged.

\subsection{Comparison with Discrete Search Baselines}

Bounded continuous optimization is strictly more effective than heuristic
discrete search for verifier manipulation. Table~\ref{tab:baseline-comparison}
compares ADSD against three discrete baselines: a pure random five-token
prefix, a GCG-adapted search baseline, and an AutoDAN-adapted search baseline.

\begin{table}[t]
\centering
\caption{Comparison with discrete search baselines on the token-wise GSM8K
benchmark. Search cost reports the dominant prompt-search workload required to
produce the final five-token prefix.}
\label{tab:baseline-comparison}
\footnotesize
\begin{tabular}{lccccc}
\toprule
Method & Search cost & Time (s) & Agg. acc. & Block eff. & GSM8K acc. \\
\midrule
Benign baseline & --- & $23.10$ & $0.597$ & $5.832$ & $222/263$ \\
Random prefix & $[XX]$ & $[XX]$ & $[XX]$ & $[XX]$ & $[XX]$ \\
GCG-adapted & $[XX]$ & $[XX]$ & $[XX]$ & $[XX]$ & $[XX]$ \\
AutoDAN-adapted & $[XX]$ & $[XX]$ & $[XX]$ & $[XX]$ & $[XX]$ \\
ADSD (UAT, $N=4$) & $[XX]$ & $38.61$ & $0.390$ & $3.731$ & $222/263$ \\
\bottomrule
\end{tabular}
\end{table}

The comparison isolates the key advantage of ADSD. Heuristic discrete methods
must search directly in token space, where verifier collapse is sparse and
projection failures are invisible until evaluation time. ADSD instead uses a
continuous verifier-aligned objective and a projection-aware regularizer to
reach deeper acceptance collapse with a short prompt budget. The resulting
prefixes therefore achieve stronger end-to-end degradation than discrete search
alone, while requiring materially less trial-and-error over the combinatorial
token space.

\subsection{Cross-Architecture Vulnerability}

The vulnerability is a property of speculative verification itself, not a quirk
of a single Qwen pairing. Table~\ref{tab:scaling-placeholder} evaluates ADSD on
larger Qwen and LLaMA draft-target pairs to test whether the same prompt-side
mechanism persists across architectures and scales.

\begin{table}[t]
\centering
\caption{Cross-architecture scaling of ADSD. Larger target models amplify the
absolute systems cost of verifier collapse because each rejected speculative
window forces more expensive target-side work.}
\label{tab:scaling-placeholder}
\footnotesize
\begin{tabular}{lcccc}
\toprule
Target $\rightarrow$ Draft & Condition & Time (s) & Block eff. & Accuracy \\
\midrule
Qwen2.5 14B $\rightarrow$ 0.5B & Benign & $23.10$ & $5.832$ & $222/263$ \\
Qwen2.5 14B $\rightarrow$ 0.5B & ADSD & $38.61$ & $3.731$ & $222/263$ \\
Qwen2.5 72B $\rightarrow$ 3B & Benign & $[XX]$ & $[XX]$ & $[XX]$ \\
Qwen2.5 72B $\rightarrow$ 3B & ADSD & $[XX]$ & $[XX]$ & $[XX]$ \\
LLaMA-3 8B $\rightarrow$ 1B & Benign & $[XX]$ & $[XX]$ & $[XX]$ \\
LLaMA-3 8B $\rightarrow$ 1B & ADSD & $[XX]$ & $[XX]$ & $[XX]$ \\
\bottomrule
\end{tabular}
\end{table}

The expected pattern is straightforward: if verifier collapse reflects a
fundamental weakness of draft-target acceleration rather than an idiosyncrasy of
one model family, then the same prompt-side mechanism should transfer across
architectures and become even more costly as target-model capacity grows. In
that regime, the absolute latency penalty becomes larger even when the relative
drop in acceptance is comparable.

\subsection{Ablation of Optimization Components}

The projection-aware design of ADSD is empirically necessary. Table
\ref{tab:ablation-placeholder} ablates the main terms of the optimization
objective, with special emphasis on the projection penalty $\rho(U)$$,$ which
is the direct algorithmic instantiation of the theoretical bound in
Section~\ref{sec:method}.

\begin{table}[t]
\centering
\caption{Ablation of ADSD optimization components. The full objective combines
early collapse pressure, trajectory-level distributional mismatch, and
projection-aware regularization.}
\label{tab:ablation-placeholder}
\footnotesize
\begin{tabular}{lcccc}
\toprule
Objective variant & Projected first token & Time (s) & Block eff. & Accuracy \\
\midrule
Full ADSD objective & $[XX]$ & $38.61$ & $3.731$ & $222/263$ \\
w/o projection penalty $\rho$ & $[XX]$ & $[XX]$ & $[XX]$ & $[XX]$ \\
w/o TV term & $[XX]$ & $[XX]$ & $[XX]$ & $[XX]$ \\
w/o reverse KL & $[XX]$ & $[XX]$ & $[XX]$ & $[XX]$ \\
w/o entropy term & $[XX]$ & $[XX]$ & $[XX]$ & $[XX]$ \\
\bottomrule
\end{tabular}
\end{table}

This ablation directly tests the central methodological claim of the paper. If
the relaxation-to-text gap were merely a nuisance detail, then removing the
projection-aware term would not materially affect the final discrete attack. Our
theory predicts the opposite: without a mechanism that keeps the relaxed prefix
close to the discrete manifold, optimization can find a strong continuous
adversarial state whose projected text no longer preserves verifier collapse.
The ablation therefore functions as an empirical validation of the
prompt-waywardness argument in Section~\ref{sec:method}.

\subsection{Discussion: Robustness Against Adaptive Defenses}

One natural defense is to detect persistent acceptance collapse and halt
drafting adaptively, as in designs such as AdaEDL. Such a fallback is useful,
but it does not eliminate the vulnerability studied here. The attacker does
not need to force the system into a permanently degraded regime to succeed; it
is already sufficient to deny the request the intended speculative speedup.
Moreover, adaptive fallbacks require an observation window before the system
can decide that drafting is no longer beneficial. The attack therefore still
extracts a penalty during the detection phase, after which the request is
reduced to ordinary autoregressive decoding rather than accelerated
speculative decoding. From the adversary's perspective, that outcome remains a
successful worst-case degradation of the algorithmic mechanism.
